% !TeX root = ../main.tex

\chapter{Background and Related Work}

\section{Price Manipulation in DeFi}

Price manipulation vulnerabilities arise when a protocol relies on a
price-related value that can be influenced by an attacker. In DeFi systems,
price is not only used for swaps. It also affects collateral valuation,
liquidation eligibility, minting and burning, vault share calculation, reward
accounting, margin requirements, and other economically sensitive state
transitions. A short-lived price deviation may therefore be enough to extract
value if the protocol reads the manipulated value and immediately commits an
irreversible state change.

Many such attacks are amplified by flash loans or other low-cost temporary
liquidity movements. The attacker first distorts the apparent market state,
then triggers a price-dependent operation, and finally restores the pool state
before the transaction ends. This makes the vulnerable condition highly
context-dependent. Detecting the issue requires understanding how a price is
obtained, how it is propagated through the program, and where it influences a
financial decision. Merely observing that a contract reads an oracle interface
or a pool state variable is not sufficient.

Common risky price dependencies include direct reads from AMM reserves, current
pool ticks, square-root prices, direct spot-price helper functions, stale
oracle answers, and protocol-specific wrappers around those values.
Common affected operations include collateral checks, liquidation logic, margin
updates, swap output calculations, liquidity minting, reward release, and
accounting state writes. A safe implementation often adds protection such as
time-weighted average prices, independent oracle cross-checks, staleness
checks, deviation bounds, slippage checks, or delayed settlement. The central
difficulty of price manipulation is therefore not recognizing one dangerous
API call, but evaluating the relationship among price source, economic sink,
and missing protection in the broader protocol context.

\section{Smart Contract Auditing Practice}

Smart contract auditing is typically conducted at the level of a complete
protocol repository rather than an isolated function or transaction. Auditors
inspect multiple contracts, helper libraries, inherited modules, deployment
assumptions, integration points, tests, and business logic together. For DeFi
protocols, this often means tracing how a price is computed, how it crosses
contract boundaries, and where it influences economically sensitive state
transitions.

The output of an audit is also richer than a vulnerability label. A typical
finding contains a title, severity, affected code location or snippet, root
cause, impact, supporting evidence or proof-of-concept context, and mitigation
recommendation. These elements matter because developers must reproduce and fix
the issue, while auditors must justify why the reported behavior is risky. An
automated tool that aims to assist audits should therefore be judged not only
by whether it predicts that a protocol is vulnerable, but by whether it
produces findings that auditors can directly inspect and validate.

\section{Related Work}

\subsection{Price Manipulation Detection Systems}

One major direction studies price manipulation through transaction-level or
post-hoc analysis. DeFiRanger reconstructs high-level DeFi semantics from raw
Ethereum transactions by building cash-flow trees and lifting low-level token
transfers into operations such as trades and liquidity changes
\citep{defiranger}. DeFiScope extends this line by combining transaction-level
DeFi operations with LLM-based abstraction of price calculation logic from
source code, and targets both standard and custom price models
\citep{defiscope}. These systems show that price manipulation cannot be
understood from syntax alone and instead requires reasoning about financial
behavior and price semantics. However, their primary output unit remains attack
or transaction detection, which is valuable for monitoring and incident
analysis but does not directly answer whether an auditor can obtain a precise
pre-deployment finding from a repository.

Another direction targets code-based or pre-mortem detection. Static analysis
is attractive in this setting because smart contract code is transparent and
many vulnerability patterns leave structural traces in control flow, data flow,
and cross-contract interactions. Slither provides a widely used foundation for
such analyses by parsing Solidity code, building intermediate
representations, and exposing program information suitable for custom security
analysis \citep{slither}. Taint analysis is especially relevant for price
manipulation because it can model whether a price-like value reaches a
financially sensitive operation. In this thesis, source and sink are analysis
abstractions rather than vulnerability labels. They identify potentially
relevant flows, while exploitability still depends on economic meaning and
missing protection.

Recent pre-mortem systems move closer to this thesis. PriceSleuth identifies
core price-calculation logic, uses an LLM to locate price-calculation
statements, and then applies dependency and propagation analysis to detect
state-manipulation vulnerabilities \citep{pricesleuth}. PMDetector combines
formal attack modeling, static taint analysis, LLM-based filtering and attack
simulation, and static validation to detect DeFi price manipulation
vulnerabilities \citep{pmdetector}. Both works are closer to audit assistance
than pure post-hoc attack detection because they reason about code before
exploitation. Nevertheless, their reported evaluation is still centered mainly
on detection performance over vulnerable and benign contracts or protocols.

Some recent systems also begin to frame the problem as an adaptive auditing
workflow. SmartAuditFlow introduces a plan-execute framework in which the LLM
dynamically revises its audit strategy based on intermediate outputs, integrates
external knowledge such as static analysis and retrieval-augmented generation,
and produces comprehensive security reports \citep{smartauditflow}. This
direction is important because it moves the field closer to multi-stage audit
reasoning rather than one-shot classification. Still, the target remains
general smart contract auditing, and the evaluation does not focus on
repository-level flash-loan price manipulation findings with explicit
source-sink-mitigation alignment.

\subsection{LLM-Assisted Detection}

Beyond price manipulation, a broader line of work studies how
LLMs can assist smart contract vulnerability detection. GPTScan is an early and
influential example that combines GPT with static analysis for smart contract
logic vulnerability detection \citep{gptscan}. Instead of treating the LLM as a
standalone oracle, GPTScan uses it as a code-understanding component to match
candidate vulnerabilities against scenario descriptions and then validates key
variables and statements through static confirmation. This hybrid design is
important because it explicitly addresses one of the main weaknesses of direct
LLM-based detection: high false-positive rates.

Other work focuses on improving reasoning structure within the LLM itself.
GPTLens frames smart contract vulnerability detection as a two-stage process in
which one LLM role generates candidate vulnerabilities and another role
discriminates among them \citep{gptlens}. The paper argues that there is a
tension between maximizing recall and controlling false positives, and that
multi-stage reasoning can partially manage this trade-off. This perspective is
relevant to audit assistance because an auditor needs both broad coverage and
low noise; simply sampling more candidate answers is not enough.

A more model-centric direction improves the smart-contract-specific capability
of LLMs through domain adaptation and explanation-oriented training.
Smart-LLaMA introduces smart-contract-specific continual pre-training and
explanation-guided fine-tuning, and explicitly targets vulnerability
detection, explanation quality, and precise vulnerability location
\citep{smartllama}. This is particularly relevant to the present thesis because
it shows that LLM-based systems can be optimized not only for binary detection,
but also for localization and explanation, which are closer to the needs of
auditing.

Another related direction treats auditing as a collaborative multi-agent
process. LLM-SmartAudit uses a multi-agent conversational framework to detect
and analyze smart contract vulnerabilities, and reports stronger performance
than traditional auditing tools on its evaluation datasets
\citep{llmsmartaudit}. While the vulnerability scope is broader than price
manipulation, the paper reflects a growing trend: LLM systems are increasingly
designed not merely to classify contracts, but to simulate aspects of audit
reasoning and report generation.

PMDetector is best understood as a price-manipulation-specific extension of
this broader LLM-assisted literature. It inherits the general move toward
hybrid reasoning, staged analysis, and richer semantic interpretation, but
applies those ideas to a narrow and economically important vulnerability class
\citep{pmdetector}. This thesis builds on the same general intuition, yet
shifts the evaluation target further toward whether the final output is useful
as a repository-level audit finding.

\subsection{Audit Report Corpora and Benchmarks}

Public audit reports and contest findings provide a natural source of real-world
security artifacts. Platforms such as Solodit aggregate findings from multiple
audit firms and contests, while Code4rena and Sherlock publish contest reports
and discussions that contain vulnerability descriptions, affected code regions,
impact statements, and remediation guidance \citep{solodit,code4rena,sherlock}.
These artifacts are valuable because they more closely resemble the output
expected from an audit assistant than coarse exploit labels or contract-level
vulnerability tags.

Recent work has begun to transform such reports into structured corpora. FORGE
uses an LLM-driven pipeline to extract self-contained vulnerability information
from real-world audit reports and classify the resulting findings under the CWE
hierarchy \citep{forge}. GiAnt likewise distills real-world audit reports into
structured findings and then applies quality assurance with an LLM-as-a-judge
mechanism to produce a large audit-oriented corpus \citep{giant}. These works
demonstrate that audit findings can be normalized into machine-usable records
with fields such as vulnerability description, severity, location, impact, and
mitigation, thereby making report-level evaluation more feasible.

In parallel, benchmark design has also shifted toward repository-level audit
tasks. ScaBench provides curated datasets from recent real-world smart contract
audits together with official scoring tools for evaluating AI audit agents
\citep{scabench}. EVMbench further expands the benchmark space by evaluating AI
agents in detect, patch, and exploit modes over curated vulnerabilities drawn
from historical audits \citep{evmbench}. Hound, in turn, uses a subset of
ScaBench to study system-level reasoning in security audits through
relation-first knowledge graphs and hypothesis-driven analysis
\citep{hound}. These benchmarks move evaluation closer to practical auditing by
treating repositories, vulnerability reports, and agent behavior as first-class
benchmark objects rather than reducing everything to file-level labels.

Taken together, the literature reviewed in this chapter establishes the main
ingredients needed by this thesis: price-manipulation reasoning, hybrid
static-analysis-plus-LLM workflows, and repository-level audit artifacts and
benchmarks. A clear gap nevertheless remains. Existing systems and benchmarks
are still largely general-purpose: they do not directly provide an evaluation
framework tailored to flash-loan price manipulation findings in which source
identification, sink localization, mitigation agreement, semantic explanation
quality, and over-reporting cost are measured together. This thesis therefore
draws methodological inspiration from these directions while focusing on a
narrower, audit-oriented evaluation target.
